\documentclass[11pt]{article}
\usepackage[margin=1in]{geometry}
\usepackage{parskip}
\usepackage{xcolor}
\usepackage{mdframed}
\usepackage{enumitem}
\usepackage{hyperref}
\hypersetup{hidelinks}

% Reviewer comment box
\newmdenv[backgroundcolor=gray!8, linecolor=gray!40, linewidth=0.5pt,
  skipabove=8pt, skipbelow=4pt, innertopmargin=6pt, innerbottommargin=6pt]{rc}
\newcommand{\resp}{\par\smallskip\noindent\textbf{Response.}\ }

\title{\vspace{-2.5em}Response to Reviewers\\[0.3em]
\large Manuscript AILET-2025-0018\\[0.2em]
\normalsize ``An Empirical Study of Deep Reinforcement Learning for Dynamic
Scheduling in Satellite QKD Networks''}
\author{}
\date{}

\begin{document}
\maketitle
\vspace{-3em}

We thank the Associate Editor and both reviewers for their careful and constructive
reviews. We have revised the manuscript thoroughly; every point is addressed below.
Reviewer comments are shown in shaded boxes, followed by our response. For
reference, the principal changes are: a fully specified reward with a provably fair
single-objective comparison, a switching-cost sensitivity analysis, a cost-aware
hybrid baseline, multi-seed (20-seed) evaluation with mean and standard deviation,
a correction to the switching-cost model, and an expanded motivation and Limitations
section. All code and models are at \url{https://github.com/ailabteam/AILET}.

\section*{Associate Editor}

\begin{rc}
We see promise in this paper, but there is still work to be done to clarify the
issues raised by the more negative of our reviewers.
\end{rc}

\resp We have focused the revision on exactly these issues. The most negative
reviewer (Reviewer 2) questioned whether the comparison was fair given an
undisclosed reward. We now specify the reward in full and make explicit that all
policies are scored by one identical physical objective (the total secure key), so
no policy can be disadvantaged by reward shaping. We also added the switching-cost
analysis and hybrid baseline requested by Reviewer 1. We believe the revised
manuscript resolves the concerns and hope it now meets the bar for publication.

\section*{Reviewer 1 (Minor Revision)}

We are grateful for the positive assessment and for recognizing the value of an
honest negative result. We address the five points in turn.

\begin{rc}
1. The evaluation is limited to a standard PPO implementation with a flat MLP
policy. The observed gap may partly reflect the choice of learning architecture and
algorithm rather than a fundamental limitation of learning-based approaches.
\end{rc}

\resp We agree and now state this explicitly in the new \emph{Limitations}
paragraph: our conclusions are specific to a single-algorithm, flat-MLP setting,
and the gap may partly reflect that architecture rather than a fundamental limit of
learning-based scheduling. We are careful throughout to frame the result as a
cautionary case study for \emph{standard} DRL, not as a refutation of learning in
general, and we point to curiosity-driven exploration, hierarchical control, and
learning--heuristic hybrids as the natural next steps. We also note that the new
per-cost ablation (point 2) shows the shortfall persists even when the agent is
optimized for each specific cost, which makes a pure tuning explanation less likely.

\begin{rc}
2. The link-switching cost is central, yet no ablation or sensitivity analysis
examines how varying its magnitude affects exploration, convergence, and final
performance.
\end{rc}

\resp We added this analysis as a core new contribution (Figure~3(b) and the
``Sensitivity to switching cost'' paragraph). We sweep the cost over
$\{0,1,2,3,4,5\}$ minutes and, importantly, retrain a DRL agent for each value.
The findings are: (i) at zero cost the DRL agent matches greedy; (ii) the gap then
widens monotonically as the cost grows, and the DRL agent degrades much faster than
greedy even when optimized for that exact cost; and (iii) beyond a critical cost of
about five minutes, once the setup latency exceeds a typical pass duration, every
policy collapses to zero because no link can be held long enough to amortize a
reconnection. This both quantifies the sensitivity and explains the mechanism behind
the original result.

\begin{rc}
3. The setup considers a single LEO satellite and five ground stations, which
restricts insight into scalability, coordination, and congestion.
\end{rc}

\resp We agree and now state this scope limitation explicitly in the
\emph{Limitations} paragraph and the Conclusion. The environment is in fact
parameterized by the number of ground stations, and we have kept the released code
configurable to support exactly the multi-station and multi-satellite studies the
reviewer envisions; we flag these as the primary directions for future work rather
than claim coverage we do not have.

\begin{rc}
4. Hybrid strategies combining learning and heuristics are mentioned conceptually
but not evaluated empirically. Including such baselines would help contextualize the
negative results.
\end{rc}

\resp We added a cost-aware hybrid baseline: greedy selection with switching
hysteresis (Eq.~9), which retains the current link unless an alternative exceeds it
in elevation by a margin large enough to justify the setup latency. To avoid a
strawman, we swept the hysteresis margin; it changes the result by at most $0.3\%$.
The hybrid does not beat plain greedy, and we explain why: with one satellite and
dispersed stations, the satellite is rarely in view of two ground stations at once,
so most reconfigurations are forced by the held link dropping below the horizon
rather than chosen to chase a marginally better one. Hysteresis can only help with
discretionary switches, which are scarce here. This contextualizes the negative
result and identifies the structural reason for it.

\begin{rc}
5. The paper would benefit from explicitly stating that these conclusions may not
directly generalize to multi-satellite, multi-agent, or long-horizon planning
scenarios.
\end{rc}

\resp Done. The \emph{Limitations} paragraph and the Conclusion now state plainly
that the conclusions need not generalize to multi-satellite constellations,
multi-agent coordination, or long-horizon planning, where learning could exhibit
clearer advantages.

\section*{Reviewer 2 (Reject)}

We thank the reviewer for the precise criticisms. We believe each rests on
information that was missing from the original submission rather than on a flaw in
the study, and we have now supplied that information.

\begin{rc}
1. There is insufficient information to be convinced. No information was provided on
how the reward function was designed, or how much weight was placed on which
elements. If the reward considers a broader range of factors than the greedy
algorithm, it will inevitably score lower on certain indicators.
\end{rc}

\resp This is the most important point and we have addressed it directly. Section~2
now specifies the complete reward as explicit equations: the free-space key-rate
model (Eqs.~4--5), the switching-cost timer dynamics (Eq.~6), and the full piecewise
reward (Eq.~7). Crucially, we clarify a point that was implicit before: \emph{the
reward is not a private DRL objective}. The single scalar in Eq.~7, the secure key
generated net of operational penalties, is the identical objective used to score
every policy, the learned agent and the greedy, hybrid, and random baselines alike.
No policy optimizes or is graded on a different or more heavily weighted quantity, so
the concern that the agent is penalized for considering ``broader factors'' does not
apply: all methods are measured on the same total secure key. We have stated this
explicitly in Section~2 and again where the results are reported.

\begin{rc}
2. The total reward in Figure 2(a) confirms convergence of the DRL algorithm. If the
rewards of the three algorithms are not designed with the same formula, it is
ambiguous what the graph explains.
\end{rc}

\resp The reviewer is right to demand this, and the answer is that the rewards
\emph{are} the same formula for all policies (see the response above). To remove any
remaining ambiguity, the learning curve (now Figure~3(a)) plots the DRL agent
against the greedy and random baselines on this one shared objective, so the curves
are directly comparable. We have also replaced the single-seed numbers with the mean
and standard deviation over 20 seeds, so the comparison reflects typical rather than
incidental behavior.

\begin{rc}
3. Why should we compare random algorithms with greedy and DRL algorithms?
\end{rc}

\resp We have clarified this in the Experimental Setup. The random policy is a
sanity-check lower bound: it confirms that the task is non-trivial and that the
reward scale is meaningful. Because the random policy earns strongly negative
returns once penalties are present (for example $-1{,}427$ in the Dynamic scenario),
it demonstrates that the constraints genuinely bite and that the positive scores of
greedy and DRL reflect real skill rather than an easy environment.

\begin{rc}
4. The title, abstract, and introduction mention the satellite QKD network, but why
the problem is treated in a network environment is not clearly described. It is hard
to understand how the experimental results affect the environment. Do you mean to
suggest that heuristics like greedy are better suited than DRL in such networks?
\end{rc}

\resp We have rewritten the motivation in the Introduction to make the network
framing concrete. A single LEO satellite is a shared downlink resource that must be
allocated, minute by minute, among several competing ground stations; the
controller's choices determine the aggregate secure key delivered across the
network, which makes this a genuine network-level scheduling problem. On the
takeaway: we are careful \emph{not} to claim that heuristics are universally better
than DRL. Our claim is narrower and, we believe, more useful: for this cost-aware
scheduling formulation, a standard DRL agent does not outperform a strong,
domain-aware heuristic, and we identify the structural reason (geometry-forced
switching) and the regime (nonzero switching cost) in which the gap appears. We
present this as a benchmark and a cautionary case for practitioners, not as a
general verdict.

\bigskip
\noindent We thank the reviewers again. We believe the revision is substantially
stronger and hope it is now suitable for ACM AI Letters.

\end{document}
